% =====================================================================
% 3. WHAT I HAD TO DO               target: 1.5 - 2 pages
%
% This section answers the examiner's second required part: "What I had
% to do". It is the MISSION, not the solution. Keep the solution for
% sections 5-8. Written mostly in the past tense.
% =====================================================================

\section{What I Had to Do}
\label{sec:objectives}

\begin{guide}
Describe the task \emph{as it was given to me}, then how I turned it into concrete
objectives. A good structure is: (1) the starting point I was handed, (2) the
problem statement, (3) the objectives with success criteria, (4) the constraints,
(5) how the plan changed along the way.
\textbf{Questions to answer:} What exactly was handed to me on day one (a paper? a
code base? nothing?). What did my supervisor ask for in one sentence? What would
``success'' have looked like? What was left open for me to decide?
\end{guide}

\subsection{Starting point}

\todo{Two or three sentences: the subject started from the \millet{} paper
\citep{early2024millet} and its public code. Say what already existed (their
backbones, their pooling heads, their metrics) and that reproducing it was step
zero, not a contribution.}

\begin{priorwork}
\todo{One or two sentences naming precisely what I reused without change: the
\millet{} training loop, the AOPCR and NDCG@$n$ metric code, the \webtraffic{} and
\ucr{} data loaders. Being explicit here protects the credibility of everything I
claim later.}
\end{priorwork}

\subsection{Problem statement}

\begin{guide}
Turn the subject into one clear research question. A good template:
``Can a much smaller encoder match the accuracy of the \millet{} baseline while
producing \emph{better} explanations?'' Then say why that question is not obvious --
smaller models usually lose accuracy, and better explanations usually cost accuracy.
\end{guide}

\todo{Three or four sentences ending with one clearly stated research question.}

\subsection{Objectives}

\begin{guide}
Three to five objectives. Each one needs a \emph{measurable} success criterion,
otherwise the conclusion cannot say whether it was reached. Come back to this list
in \Cref{sec:conclusion} and answer each point honestly.
\end{guide}

\begin{table}[h]
\caption{The objectives I set with my supervisor, and how each one was measured.
\todo{Fill in the criteria; keep them measurable.}}
\label{tab:objectives}
\begin{center}
\small
\begin{tabular}{p{0.32\linewidth} p{0.36\linewidth} p{0.22\linewidth}}
\toprule
\textbf{Objective} & \textbf{Success criterion} & \textbf{Reached?} \\
\midrule
O1. \todo{Reproduce the baseline} & \todo{e.g. accuracy within X of the published
number under my own training recipe} & \tbd \\
O2. \todo{Design a smaller encoder} & \todo{e.g. $\geq$ 50\,\% fewer parameters at
equal accuracy} & \tbd \\
O3. \todo{Improve explanation quality} & \todo{e.g. higher AOPCR and NDCG@$n$ on
\webtraffic{}} & \tbd \\
O4. \todo{Build a reproducible benchmark} & \todo{e.g. any model reproducible from
one command and one config file} & \tbd \\
O5. \todo{optional fifth objective} & \todo{criterion} & \tbd \\
\bottomrule
\end{tabular}
\end{center}
\end{table}

\subsection{Constraints}

\begin{guide}
Short bullet list. Constraints make the work look real, and they explain later
choices (for example: why only three seeds, why \webtraffic{} was the screening
dataset).
\textbf{Questions to answer:} How much time did I have? How much GPU time? Which
comparisons had to stay fair (same recipe for every model)? What could I not change
(the metric definitions, the data splits)?
\end{guide}

\begin{itemize}\itemsep2pt
  \item \todo{Time: \num{months} months of internship.}
  \item \todo{Compute: shared GPU nodes with a job time limit; a full sweep over 128
        datasets takes \num{hours} hours per model.}
  \item \todo{Fairness: every model trained with the same recipe, so that only the
        architecture differs.}
  \item \todo{Comparability: keep \millet{}'s metric code untouched, so my numbers
        can be read next to theirs.}
\end{itemize}

\subsection{How the plan evolved}

\begin{guide}
Very valuable for the examiner: it shows scientific judgement rather than blind
execution. Two or three sentences. Be honest -- a plan that changed for a good
reason is a strength, not a weakness.
\textbf{Questions to answer:} What did I try first that did not work? What made me
change direction (a result? a paper I read? a supervisor's remark?). Where did the
idea of borrowing pooling ideas from cancer-pathology MIL papers come from?
\end{guide}

\todo{Three or four sentences describing the path: baseline reproduction $\rightarrow$
smaller encoder $\rightarrow$ pooling heads $\rightarrow$ ideas imported from
whole-slide-image MIL \citep{ilse2018attention, li2021dsmil} $\rightarrow$ multi-seed
validation.}

% Optional but nice: a small timeline figure of the internship.
\begin{figure}[h]
  \centering
  \figph[0.85]{Timeline of the internship: reading and reproduction, encoder design,
  pooling design, large-scale sweeps on Grid'5000, analysis and writing. One
  horizontal bar per phase, weeks on the $x$ axis.}
  \caption{\todo{Caption: how the \num{months} months were spent.}}
  \label{fig:timeline}
\end{figure}
